\documentclass[10pt,a4paper]{amsart}
\usepackage[margin=19mm]{geometry}
\usepackage{amsmath,amssymb,mathtools}
\usepackage[scr=rsfs,cal=euler]{mathalfa}
\usepackage{xfrac}
\usepackage[unicode,hidelinks,bookmarks=false]{hyperref}
\newcommand{\ie}{i.e.\ }
\newcommand{\cf}{cf.\ }
\newcommand{\resp}{respectively\ }
\newcommand{\Subsection}{\subsection}
\providecommand{\nobreakdash}{\nobreak\hbox{-}}
\setlength{\emergencystretch}{2em}
\multlinegap0pt
\usepackage{bm}
\usepackage{bbm}
\usepackage{dsfont}
\usepackage{xspace}
\usepackage{cancel}
\usepackage{mathtools}
\usepackage{amsthm}
\makeatletter
\@ifundefined{theorem}{\newtheorem{theorem}{Theorem}}{}
\@ifundefined{proposition}{\newtheorem{proposition}[theorem]{Proposition}}{}
\makeatother
\renewcommand{\Re}{\mathsf{Re}\,}        
\title[Spectrum of the wave equation with Dirac damping on a compac]{Spectrum of the wave equation with Dirac damping on a compact star graph}
\author{Mikul\'a\v{s} Ku\v{c}era}
\date{Source: arXiv:2510.09829v2 (CC BY 4.0)}
\begin{document}
\maketitle

\noindent\emph{Source and licence.} Spectrum of the wave equation with Dirac damping on a compact star graph. Version arXiv:2510.09829v2 (CC BY 4.0), distributed under the Creative Commons Attribution 4.0 International licence, as recorded in the arXiv licence metadata for this exact version and shown on the abstract page https://arxiv.org/abs/2510.09829v2. Full rights notice: licenses/arxiv-2510.09829-CC-BY-4.0.txt.\par\smallskip

\noindent\textit{Editorial scope.} Counted unit: the Proposition whose source label is \texttt{P: Roots at most double} in the article (source0.tex lines 594--596, source label \texttt{P: Roots at most double}), delivered together with the article's own complete original proof of it (source0.tex lines 597--655, one \texttt{proof} environment of 59 lines). No mathematics is written, repaired or re-derived: statement and proof are the article's own text line for line. Required context retained uncounted: none. Every definition, display and label of those lines is retained unchanged. Omitted from this package and not redistributed: 1--593, 656--1818. Nothing inside a retained excerpt is omitted or altered: every retained body is delivered exactly as the article's lines are written, so no omission receipt of any kind accompanies this package. Citations: the retained excerpts carry no citation command at all, so no bibliography entry of the article is transcribed and no reference is redistributed. Reference adaptation: none was needed and none was made; every reference inside the delivered unit resolves inside it, so no unresolved reference or citation is produced. Printed slips kept literal: no printed word, symbol, hypothesis or number of the counted statement or of its proof was altered. Layout and class: the article's own class is \texttt{amsart} with packages including subcaption, color, hyperref, babel, amsmath, amssymb, amsthm, graphicx, amsfonts, url, enumerate, dsfont, stmaryrd, mathabx; the wrapper is the lane's pinned \texttt{amsart} editorial wrapper at 10pt on A4, which restates the theorem-like environments the retained text needs and adds the article's own macro definitions that the retained text needs (\texttt{\textbackslash Re}), with extra wrapper packages (bm, bbm, dsfont, xspace, cancel, mathtools). The delivered numbering is the unit's own. Assets: no retained span contains a figure environment, a graphics-inclusion command, a PDF-page inclusion, a table or a TikZ picture; no image file is copied into this package, the package declares no asset dependency and no image file is redistributed with it. Licence: version arXiv:2510.09829v2 of this article is distributed under the Creative Commons Attribution 4.0 International licence, as recorded in the arXiv licence metadata for this exact version and shown on the abstract page https://arxiv.org/abs/2510.09829v2. A full rights notice is delivered at licenses/arxiv-2510.09829-CC-BY-4.0.txt. Characters: every retained excerpt is pure ASCII, so the delivered engine, XeLaTeX with its default Latin Modern text font, reports no missing character.
\begin{proposition}\label{P: Roots at most double}
    All roots of $S(\cdot; a, \alpha)$ are at most double.
\end{proposition}

\begin{proof}
    Let $F(\lambda) \coloneqq \lambda S(\lambda; a, \alpha)$ for simplicity. Note that
    \[
    F(\lambda) = \sinh(\lambda\pi) + \frac{\alpha}{2}\cosh(\lambda\pi)-\frac{\alpha}{2}\cosh(\lambda(\pi-2a)).
    \]

    We have
    \begin{equation}
    F'(\lambda) = \pi\cosh(\lambda\pi) + \pi\frac{\alpha}{2}\sinh(\lambda\pi) - (\pi-2a)\frac{\alpha}{2}\sinh(\lambda(\pi-2a))
    \label{FDerivative}
    \end{equation}
    and for the second derivative
    \begin{align*}
    F''(\lambda) &= \pi^2\sinh(\lambda\pi) + \pi^2 \frac{\alpha}{2}\cosh(\lambda\pi) -(\pi-2a)^2\frac{\alpha}{2}\cosh(\lambda(\pi-2a)) \\
    &= \pi^2F(\lambda) + 2a\alpha(\pi-a)\cosh(\lambda(\pi-2a)).
    \end{align*}

    Suppose for a contradiction that 
    \begin{equation}
    F(\lambda_0) = F'(\lambda_0) = F''(\lambda_0) = 0
    \label{eq: Triple root}
    \end{equation}
    for some $\lambda_0 \in \mathbb C$. Let us first treat the case $\alpha = 0$. Then $F(\lambda_0) = \sinh(\lambda_0\pi) = 0$ means that $F'(\lambda_0) = \pi\cosh(\lambda_0\pi) \neq 0$ which contradicts \eqref{eq: Triple root}.
    
    
    If $\alpha \neq 0$, then \eqref{eq: Triple root} forces 
    \[
    \cosh(\lambda_0(\pi-2a)) = 0 \implies \lambda_0 = \frac{\mathrm{i}\pi\left(n + \frac{1}{2}\right)}{\pi - 2a}, \quad n \in \mathbb Z.
    \]
    The implication is only valid if $a \neq \pi/2$. In the case $a=\pi/2$, the left-hand side becomes $1=0$, which is a contradiction.


    Simultaneously,
    \begin{equation}
    0 = F(\lambda_0) = \sinh(\lambda_0\pi) + \frac{\alpha}{2}\cosh(\lambda_0\pi) = \sinh(\mathrm{i}\omega) + \frac{\alpha}{2}\cosh(\mathrm{i}\omega) = \mathrm{i}\sin\omega + \frac{\alpha}{2}\cos\omega,
    \label{DoubleEigenvaluesProof}
    \end{equation}
    where we denoted $\omega \coloneqq \frac{\pi^2\left(n + \frac{1}{2}\right)}{\pi-2a}$.

    If $\Re \alpha \neq 0$, then $\cos\omega = 0$ and consequently also $\sin\omega = 0$ -- a contradiction.
    
    If $\alpha \in \mathrm{i}\mathbb R$, the equation \eqref{FDerivative} for $F'(\lambda) = 0$ yields $\cos\omega = \cosh(\lambda_0\pi) = 0$ by taking the real part. However, \eqref{DoubleEigenvaluesProof} then forces also $\sin\omega = 0$ giving us the same contradiction as above.
    % However,  the equation for $F'(\lambda_0) = 0$ yields
    % \[
    % \cos\omega + \mathrm{i}\frac{\alpha}{2}\sin\omega = 0,
    % \]
    % giving us
    % \[
    % \begin{pmatrix}
    %     \mathrm{i} & \frac{\alpha}{2} \\
    %     \mathrm{i}\frac{\alpha}{2} & 1
    % \end{pmatrix}\begin{pmatrix}
    %     \sin\omega \\ \cos \omega
    % \end{pmatrix} = \begin{pmatrix}
    %     0 \\ 0
    % \end{pmatrix}.
    % \]
    % Note that the determinant of the matrix is non-zero if and only if $\alpha \neq \pm 2$. In such case, both $\sin\omega$ and $\cos\omega$ must be zero -- a contradiction. If $\alpha = \pm 2$, then \eqref{DoubleEigenvaluesProof} is sufficient to get the same outcome.
\end{proof}

\end{document}
